% !Mode:: "TeX:UTF-8"
\chapter{燕山大学研究生学位论文撰写规范}
\label{appendix}
% 重新定义 公式编号格式、图形、表格
\renewcommand\theequation{A-\arabic{equation}}
\renewcommand\thefigure{A-\arabic{figure}}
\renewcommand\thetable{A-\arabic{table}}


\begin{equation}\label{A1}
a=b^2+c_2
\end{equation}

学位论文是研究生科学研究工作的全面总结，是表述其研究成果、代表其研究水平的重要学术文献资料，是申请和授予相应学位的基本依据。撰写学位论文是研究生培养的基本训练之一，必须按照规范认真执行。指导教师应加强指导，严格把关。
\begin{figure}[htbp]
\centering
\includegraphics[width=8cm]{ysulogo}
\caption{The Yanshan University LOGO}
\end{figure}\\
撰写的论文应符合国家及专业部门制定的有关标准，符合汉语语法规范。

硕士和博士学位论文除学术水平不同外，它们的撰写要求基本一致。

\begin{table}[htbp]
 \centering\zihao{5}
 \caption{燕山大学硕士学位论文参考文献规则}\label{tab:ysubof1}
 \begin{tabular}{llr}
 \toprule
    论文版本    & 参考文献标准    & 实施年份（年）  \\
 \midrule
    旧版        & BF7714-87       & 1987            \\
    新版        & GBT7714-2005    & 2005            \\
 \bottomrule
 \end{tabular}
\end{table}

\section{封面（封一、封皮）}\label{appendixA-1}

参见附录1学位论文封面示例。

封面格式参照附录1。各项名称的横向位置应居中安排，竖向间距按附录1安排。

开本尺寸，博士论文的为A4，装订切齐后为205×290mm；硕士论文的为B5，装订切齐后为178×250，单位为mm。

博士论文的书脊处应印：论文题目（小3号黑体，距上边线50mm）、姓名（距下边线30 mm）。

\section{封里（封二）}\label{appendixA-2}

博士、硕士学位论文的封里包括中文页及英文页。中、英文页的版式相同。对于博士、硕士等不同情况者的学位论文封里可参考附录2。

以研究生毕业同等学力申请学位者，需在封里的右上角加上如下字样：（同等学力人员）。

\section{摘要及关键词}\label{appendixA-3}

摘要是学位论文的必要附加部分，是以提供学位论文内容梗概为目的，不加评论和补充解释，简明、确切记述学位论文重要内容的完整短文。它反映了学位论文的精华。摘要应以浓缩的形式概括学位论文研究工作的目的、课题、基本观点、主要内容、研究方法及结论等。摘要应具有独立性和自明性，并拥有与学位论文同等量的主要信息，即不阅读学位论文全文就能获得必要的信息。

博士及硕士学位论文的摘要均要求用中、英两种文字给出，中文在前。

摘要的字数，以汉字计，硕士学位论文为500～650字；博士学位论文为900～1200字。摘要页不写论文题目。摘要中不宜使用公式、图表及参考文献序号。

摘要页的其它要求见附录3摘要及关键词示例。

关键词是为了适应计算机自动检索的需要，从论文中选取的能够反映文献特征内容的规范性词（称叙词或主题词）。应按GB3860-83《文献主题标引规则》的规定选取。对于那些反映新学科、新技术中的重要术语，也可作为关键词标出。关键词一般列5～10个，按词条外延层次从大到小排列。

\section{学位论文的组成部分和排列顺序}\label{appendixA-4}

学位论文一般由以下几个部分组成：封面、论文摘要、论文目录、正文、参考文献、
发表文章目录、致谢等。

\section{目录}\label{appendixA-5}

博士学位论文目录用中、英文给出，硕士学位论文目录只用中文给出。目录页的其它要求见附录4目录示例。

\section{正文}\label{appendixA-6}

\subsection{正文的内容}\label{appendixA-6-1}

正文由绪论、本论和结论三部分组成。

\subsubsection{绪论}\label{appendixA-6-1-1}

绪论是学位论文主体部分的开端，是全篇论文的引子。绪论主要回答``为什么研究（Why）"这个问题。它主要包括：

(1)研究的出发点与目的；

(2)研究的主要问题及范围，即研究主题简介；

(3)与学位论文有关的重要文献资料综述，相关领域前人及他人研究成果、水平和知识空白，存在的主要问题以及急待解决和完善的问题；

(4)理论分析依据，研究设想，研究方法和实验设计的概述；

(5)预期研究结果，科学意义和实用价值。

为了说明学位论文作者对研究方案的充分论证，反映作者在相关的学科领域中达到的学识水平以及具有开阔的科学视野，通常将文献综述作为绪论的重要内容。

综述用大量篇幅对本研究课题进行历史回顾，对前人和他人的研究工作及成果进行综合评述，然后引出自己的课题研究内容。综述相当于在学位论文研究领域内进行一次信息资料的专题研究报告。它反映了研究生阅读有关著作、学术期刊和文献资料的情况；也反应他所具备的基础和专业知识以及独立的自学能力。综述要求研究生对所阅读的文献资料进行消化、分析、判断和综合，决定取舍，确定自己所应开展科研工作的内容。只有在认真细致查阅参考文献的基础上，才能使自己的研究选题具有可靠的基础，尽量避免重复前人的研究工作，从而在前人的基础上解决前人没有解决或有待完善的工作。

撰写学位论文的绪论时，不要与摘要雷同，不要诠释基本理论，不要推导基本公式，不要介绍研究方法细节。写法应开门见山，言简意赅，措词精炼。

\subsubsection{本论}\label{appendixA-6-1-2}
本论是论文的主体部分，是作者研究成果的展示和表述。主要回答``怎么研究（how）"这个问题。它针对绪论部分提出的论点，全面透彻的进行分析论证，充分表达作者的见解。中心论点的论证可以从不同的角度展开，因此，可建立不同的分论点，共同围绕中心论点为其服务。

本论在安排层次上主要有三种方法：①递进式，即先提出一个论点，然后步步深入，层层展开，循序进行论述；②并列式，即把中心论点的几个从属论点并列起来，分别加以论述。③结合式，即把递进式和并列式结合在一起使用，这种方式多用于内容较多、篇幅较长的学位论文。层次可用小标题或序码表示。

本论的论证方法主要有：归纳分析法、演绎分析法、例证法、引证法、动因分析法、比较分析法、因素分析法、综合评析法等，可视情况灵活选用。

本论通常占学位论文的大部分篇幅。它的具体陈述方式往往不同学科有较大的差别，不能牵强地作出统一的规定。一般应在论文总体框架下按``概念的提出，理论的论证，方法的选择，试验的观察，数据的处理与分析，结果的检验"等，逐层分析，分章撰写。每一章的最后都要有``本章小结"。本章小结是用实践和数据给出的分析研究结果，最忌主观推测成分。
本
论要求思路清晰，合乎逻辑，用语简洁准确、明快流畅；内容务求客观、科学、完备，要尽量让事实和数据说话。凡是用简要的文字能够讲清楚的内容，应用文字陈述。用文字不容易说明白或说起来比较繁琐的，应用表或图来陈述。数据的引用要严谨确切，资料的引用要标明出处。

教科书式的撰写方法是撰写学位论文的第一大忌。对已有的知识避免重新描述和论证，尽量采用标注参考文献的方法；对用到的某些数学辅佐手段，应防止过分注意细节的数学推演，需要时可采用附录的形式收入学位论文供读者选阅。

\subsubsection{结论}\label{appendixA-6-1-3} 
结论是指学位论文最终的总体结论。主要回答``研究出什么（what）"这个问题。它是在对本论分析、综合的基础上对全文的总结和升华，是对绪论和本论论点的呼应，是对全文的画龙点睛的收束，也是对本论中各章小结的再综合。可以作为结论的主要内容有：

（1）	本课题的理论、方法、设计计算和试验研究结果等的归纳和总结；

（2）	明确回答本研究课题所提出的问题和任务解决的程度和取得的进展；

（3）	明确指出本研究内容的创造性成果或创新点理论（含新观点、新思路、新见解）和方法；

（4）	对研究成果应用前景和社会、经济价值等加以预测和评价；

（5）	指出今后在本研究方向进一步开展研究工作的设想、展望和建议，仪器设备的改进，尚待解决的问题。

撰写结论时应注意的事项：

（1）	结论应该明确、精炼、完整、准确、措辞严密，不含糊其词；

（2）	结论要一分为二，要尊重前人的研究成果，实事求是的评价自己的成果，切勿言过其实，在无充分的把握时，应该留有余地。

（3）	结论应反映自己从事研究工作取得的成绩，属于前人或他人已有的正确结论，应不提或尽量少提。

结论应分条撰写，字数一般不多于2000字。


\subsection{正文的撰写要求}\label{appendixA-6-2}

\subsubsection{学位论文字数}\label{appendixA-6-2-1} 
博士学位论文正文字数，理工科为8～12万字；管理及人文学科为10～14万字，其中绪论占一万字左右。

硕士学位论文正文字数，理工科为4～6万字，管理学科5～7万字（MBA3～4万字，案例分析论文不少于2万字），人文学科为3～5万字，外国语学科不少于2万外文词。其中绪论（人文学科为``导言"，外国语学科为``引言Introduction"）占3000～5000字。

\subsubsection{学位论文的版面要求}\label{appendixA-6-2-2}

(1) 页眉

页眉应居中置于页面上部。单数页眉的文字为``章及标题"；双数页页眉的文字为``燕山大学硕士（博士）学位论文"。页眉的文字用5号宋体（文字上空一行书写）。页眉文字下面为两条横线(两条横线的长度与版芯尺寸相同，硕士论文线粗0. 5磅，博士论文线粗1.5磅）。

(2) 页码

摘要和目录的页码用罗马数字，正文的页码用阿拉伯数字，居中标于页面底部。正文部分的首页和翻开后的每一右页都应该是单数页码。研究生学位论文的原件要求用计算机排版、编辑与打印，余者可复印。

(3) 正文层次编排

建议采用表1所示格式，但若节下面无需列条者也可直接列款或项。

章标题用小2号黑体字，横向居中排放，竖向位置与目录页页题上、下留空相同，题序与题名间留一字空。

节标题用小3号黑体，条标题用4号黑体，款标题用小4号黑体，正文用小4号宋体字。

外国语学科论文应使用英文Times New Roman小4号体，不可使用花体字；书名、文献名等作品名称应使用斜体字，段落书写采用首行缩进格式。

各层次标题均不得置于页面的最后一行，即不允许``背题"。

博士学位论文开本为A4，页边距为设置：上下分别为3cm，左右分别为2.8cm；正文满页为30行，每行36个字，每页版面字数为1080，行间距为最小值22磅。硕士学位论文开本为B5，页边距为设置：上下分别为2.7cm和2.3cm，左右分别为2.2cm和2.2cm，正文满页为29行，每行32个字，每页版面字数为928，行间距为最小值20磅。

\subsubsection{引用文献}\label{appendixA-6-2-3} 
正文中引用文献，应将所引用文献的编号（用阿拉伯数字的小4号字）括以方括号，置于所引内容最后一个字的右上角，如``二次铣削[1]"。当引用某文件的结论时，可将该文件编号用小4号字加方括号与正文排齐，如``由文献[8,10～14]可知"。

各级标题处均不得设置引用文献标示。

\subsubsection{名词术语}\label{appendixA-6-2-4} 
科技名词术语及设备、元器件名称，应符合有关标准中的规定。标准中未予规定的术语、名词，应采用行业通用的。论文中自行采用的新名词（包括中、外文的）应加以注释。第一次出现的英文缩写应加括号标出全称。

\subsubsection{物理量的表示}\label{appendixA-6-2-5} 
论文中表示物理量的名称和符号应执行GB3100~3102-86的规定（参见附录5）。物理量的符号都必须用斜体。

\subsubsection{计量单位}\label{appendixA-6-2-6} 
物理量计量单位及符号应按国务院1984年发布的《中华人民共和国法定计量单位》及GB3100～3102《量和单位》系列国家标准（见附录5、6）执行，不得使用非法定计量单位及符号。

表达量时，在公式图表和文字叙述中，一律使用单位的国际符号。某些单位没有国际符号时，可用汉字与国际符号构成组合单位。如：m2/人，t/月。数值与单位符号间应留一个空格。
量符号必须采用斜体字母。矢量、张量符号一律用黑斜体。量符号下角标有斜体的，正体的和正体斜体混合的，要注意区分（见附录7）。

计量单位符号一律用正体。一般单位符号为小写体，只有来源于人名的单位，其符号的首字母大写。

词头是为了避免过大或过小的数值而加在SI单位之前构成十进倍数或分数单位的因数符号。

词头符号一律采用正体。表示因数符号大于106的词头符号为大写体，其余均为小写体。特别要注意区分Y（1024）和y（10－24）、Z（1021）和z（10－21）、P（1015）和p（10－12）、M（106）和m（10－3）。词头符号与单位符号之间不留空隙，也不加任何符号。

\subsubsection{数字}\label{appendixA-6-2-7} 
按国家语言文字工作委员会等七单位1987年公布的《关于出版物上数字用法的执行规定》，除习惯用中文数字表示的以外，一般均采用阿拉伯数字（见附录8）。

\subsubsection{数理公式}\label{appendixA-6-2-8} 
公式原则上应居中书写。比较简单的或叙述性的式子可以串文排。

公式序号按章编排，如第2章第1个公式为``（2-1）"。序号加圆括号，与公式同行排在右顶格处。

文中引用公式时，用``见式（1-1）或由公式（1-1）"。公式后若有说明文字，如``式中……"，应另起一行，左起顶格排放。

为了节省版面，便于转行或减少排版难度，要对公式采用合理的表达形式，如 和 ， 和 ， 和 等，采用后者较方便。
公式转行容易出错。需要转行时，最好参阅一下数理公式的转行规则（见附录9）。

\subsubsection{插表}\label{appendixA-6-2-9} 
学位论文中表格的运用是比较多的。表格能够增强表述内容的逻辑性和准确性，已成为现代科技文献中不可缺少的表述手段。使用合适的表格，将使论文的篇幅紧凑、论述清晰，给人以强烈的对比效果。
一个表格应尽量保持完整，没有特殊需要不要分割成两部分或更多部分。表格的位置安排，一般应随文列出，要紧接在第一次涉及它的文字段后面，应尽量与涉及文字在同一段落，或编排在同一页码上。但有时限于论文结构，或表格容量较大，也可以把表格放在涉及文字段落稍后的地方，这时仍然应该争取与涉及的文字段落放在同一章节和同一个视页中。必要时，表格也可分为两段或多段（这只能发生在转页时），转页分段后的每一续表的表头都应重新排字，重排表头的续表起始横线上方居中应注明``续表 "字样。

有些表宽大于版芯的表格，需要左转90o排，这时表头一律朝左（即单页上朝钉口，双页上朝切口），表底一律朝右。少数容量特别大的表格需要跨页排时，应尽量采用双页跨单页底骑缝表方式（避免单跨双的转面表），以使整个表格呈现在一个视野上。
表格不加左右边线，上、下边线为粗实线（1.5磅）。

表格均应有表题（由表序和表名组成）。表序编排规则与公式符号的一样。表序与表名之间空一格，表名中和表名后均不使用标点符号。表题置于表上，居中排放。表题要求用中、英两种文字，中文在上，中、英文表题以及表题与表格间距为最小值18磅。工程硕士学位论文不要求英文表题。
全表如用同一单位，应将单位符号移至表头右上角，加圆括号(见附录10中例2）。

表内数字空缺的格内加横线``-"（占两个数字宽度）。表内文字或数字上、下或左、右相同时，采用通栏处理方式（见附录10中例2），不允许使用``〃"、``同上"之类的写法。

表内文字说明，起行空一格，转行顶格，句末不加标点。表题用5号字，表内文字及表的说明文字均用小5号字，英文用正体，中文用宋体。

\subsubsection{插图}\label{appendixA-6-2-10} 
学位论文中的插图与文字一样是用来表达作者意图的有效工具。插图可以形象、直观地表述科学思想和技术内容，被称为``形象语言"，插图可以使某些内容的描述简洁、清晰，可省略大段文字，有利于紧缩篇幅。
插图重要特点是它的原始设计，应由论文作者自己来完成。设计插图时应注意图形的示意性；内容的写实性；取舍的灵活性和绘制的规范性。插图的线条粗细与插图的大小有关，插图大线条粗些，插图小线条细些，要保持整篇论文插图线条的统一。
论文中使用的机械工程图应符合GB4457～GB131-83《机械制图》标准规定；电气图应符合附录11所列有关标准的规定；流程图原则上应采用结构化程序并正确运用流程框图。对于无规定符号的图，应采用相应行业上的常用画法。
每个图均应有图题（由图号和图名组成）。图号编排同表序编排规定。图题置于图下，表题用中、英文两种文字居中排放。中文的在上。文字与图的竖向间距同插表。图名在图号之后空一字排放。引用图要在图题右上角加引用的文献号。图中若有分图时，分图图题也置于分图之下，分图号用(a)、(b)等表示。分图应尽量横向排列，分图较宽者，可纵向排列，分图较多者，也可横、纵向组合排列。图宽大于版芯的图，需要左转90o排，单数页图题在切口边，双数页图题在钉口边。
图题用5号字，图内文字用小5号字。中文为宋体，英文为正体。工程硕士学位论文不要求英文图题。

论文中引入插图之前，必须有该插图的提示语，如``见图1-1"等。若提示语后面的版页不够排放该图时，则可将图后部分文字提到图前排版，将该图移到次页最上端或次页首段末。


\subsubsection{照片图}\label{appendixA-6-2-11} 
学位论文原件中照片图应采用原版照片（黑白、彩色均可）粘贴，不得用复印图。照片采用光面相纸不用布纹相纸。对金相显微组织照片必须注明放大倍数。

学位论文原件中的插图也不得用复印件。对于复杂的引用图，可采用数字化仪输入计算机打印出来的图。

\section{附录}\label{appendixA-7}

附录是学位论文的附件，不是必要组成部分。对需要收录于学位论文中且又不适
合置于正文的附加数据、资料、公式推导，程序等内容，可做为附录排在``结论"和``参考文献"之间，序号采用``附录1"、``附录2"等，采用小2号黑体字，左侧顶格开排，纵向位置同``目录"

\section{参考文献}\label{appendixA-8}

在学位论文中收录参考文献具有多方面的作用：1）对他人研究成果的尊重；2）
提供阅读学位论文的参考资料；3）表明学位论文中所引资料的可信度；4）便于他人参阅从事进一步研究。

（1）	准确恰当地引用和标注参考文献有助于提高学位论文的学术水平。
 撰写学位论文都必然涉及到对本领域前人或他人研究成果的比较、分析、综合
或者运用，以便用来支持和丰富自己的学位论文。准确恰当地引用和标注参考文献，有助于反映学位论文的学术背景、研究目的和研究深度；有助于反映学位论文作者的科学态度和科学依据；有助于说明研究方法或过程的合理性以及研究结果的可靠性；有助于明确本人研究与前人研究的关系，突出学位论文的创新或独到之处；有助于明晰学位论文和所引用文献的著作权关系，避免因为漏引、误引或标注不清所引起的抄袭或者剽窃之嫌。

（2）	准确恰当地引用和标注参考文献是科学研究工作者的职业道德。
 既是对前人或他人劳动成果的尊重，又是对提高学位论文的学术质量有水涨船
高之效。

（3）	准确恰当地引用和标注参考文献是对参考文献作者的尊重。
故意把他人作品中的成果（说法、公式、数据、结论等）当作自己的而不加以说
明，视为抄袭或剽窃行为。

（4）	准确恰当地引用和标注参考文献是培养研究生树立知识产权意识的一项技
术性很强的工作。

    参考文献书写格式应符合BF7714-87《文后参考文献著录规则》。常用参考文献编写项目和顺序规定如下：

          著作图书文献
  序号    作者.书名.版次.出版地:出版者,出版年:引用部分起止页
                                                         第一版应省略
          翻译图书文献
  序号    作者.书名.译者.版次.出版地:出版者,出版年:引用部分起止页
                                                         第一版应省略
          学术刊物文献
  序号    作者.文章名.学术刊物名,年,卷(期):引用部分起止页
          学术会议文献
  序号    作者.文章名.编者名.会议名称,会议地址,年份.出版地,出版者,出版年:引用         部分起止页
           学位论文类参考文献
  序号    研究生姓名.学位论文题目.[学校及学位论文级别].答辩年份:引用部分起止页（起止页码之间采用短横线``-"表示）

    学术会议若出版论文集者，可在会议名称后加上``论文集"字样。未出版论文集者省去``出版者"、``出版年"两项。会议地址与出版地相同者省略``出版地"。会议年份与出版年相同者省略``出版年"。
参考文献字号用5号宋体。参考文献书写格式见附录12参考文献示例。
博士学位论文的参考文献数一般应不少于100篇，其中外文文献一般不少于
总数的1/2；硕士学位论文的参考文献一般应不少于50篇，其中外文文献一般不少于20篇。参考文献中近五年的文献数一般应不少于总数的1/3，并应有近两年的参考文献。工程硕士学位论文的参考文献一般不应少于20篇。
人文学科的``参考文献部分"分为``注释"和``参考文献"两部分。其中"注释"采用脚注的形式。注释和参考文献总数要求在50篇以上，其中注释大于30篇，注释和参考文献中的外文文献不少于1/10。
外国语学科学位论文注释采用尾注的形式。

    产品说明书、各类标准、各种报纸上刊登的文章及未公开发表的研究报告（著名的内部报告如PB、AD报告及著名大公司的企业技术报告等除外）等通常不宜作为参考文献引用。
引用网上参考文献时，应注明该文献的准确网页地址，网上参考文献不包含在上述规定的文献数量之内。


\section{攻读学位期间承担的科研任务与主要成果}\label{appendixA-9}

该页页题为：攻读博（硕）士学位期间承担的科研任务与主要成果，其字体、字
号及横、纵向位置同参考文献页。该项的主要内容包括：

(1)	承担的研究课题；

(2)	科学研究成果（包括已发表的学术论文；研究报告；专著与教材；科研成果获奖项目等）；

(3)	专利申报情况。

\section{致谢}\label{appendixA-10}

论文的最后一段写致谢，这是作者对论文的形成作过贡献的组织或个人予以感谢
的文字记载。致谢的对象主要有：

1.	对资助研究工作的基金会、合作单位、企业和个人表示谢意；

2.	对指导教师、答疑教师、实验教师以及协助完成研究工作的人员致谢；

3.	对论文撰写过程中提出建议和提供帮助的人员致谢；

4.	对学位论文引用过的资料、图表、文献、研究思想和设想的知识产权所有者表示感谢。

致谢不仅是一种礼貌，也是对他人劳动的尊重。致谢的内容要实在，语言要诚恳、恰当、简练；不应该借学位论文完成之机，对提供过与论文工作无关的其它帮助的人员致谢。

\section{索引}\label{appendixA-11}

为便于索引文中内容，可在论文``致谢"页后编排索引页（若无需要亦可不设索引页）。索引以论文中的专业词语为检索线索，指出其相关内容所在页码。索引设中、英两种文字页，中文页在前，中文页按名词汉语拼音第一个字母排序，英文按该词第一个英文字母排序。索引示例见附录13。

\section{作者简介}\label{appendixA-12}

所有学科的博（硕）士学位论文都要求附以``作者简介"页。左上角放置作者彩照（小2寸）。

作者简介包括的内容：姓名、性别、民族、籍贯。作者简历：大学毕业的时间、院校、专业、学位。主要工作单位及从事过的与本学位论文有关的工作。取得的科研成果：发表的学术论文、专著、教材。获奖情况及获国家专利情况等。
